
\documentclass{article}
\usepackage{colortbl}
\usepackage{makecell}
\usepackage{multirow}
\usepackage{supertabular}

\begin{document}

\newcounter{utterance}

\centering \large Interaction Transcript for game `taboo', experiment `high\_en', episode 15 with reka{-}flash{-}3.1{-}t0.0.
\vspace{24pt}

{ \footnotesize  \setcounter{utterance}{1}
\setlength{\tabcolsep}{0pt}
\begin{supertabular}{c@{$\;$}|p{.15\linewidth}@{}p{.15\linewidth}p{.15\linewidth}p{.15\linewidth}p{.15\linewidth}p{.15\linewidth}}
    \# & $\;$A & \multicolumn{4}{c}{Game Master} & $\;\:$B\\
    \hline

    \theutterance \stepcounter{utterance}  
    & & \multicolumn{4}{p{0.6\linewidth}}{
        \cellcolor[rgb]{0.9,0.9,0.9}{
            \makecell[{{p{\linewidth}}}]{
                \texttt{\tiny{[P1$\langle$GM]}}
                \texttt{You are playing a collaborative word guessing game in which you have to describe a target word for another player to guess.} \\
\\ 
\texttt{Rules:} \\
\texttt{(a) You have to reply in the form: CLUE: <some text>. Guesses from the other player will start with GUESS.} \\
\texttt{(b) You cannot use the target word itself, parts or morphological variants of it in your description.} \\
\texttt{(c) In addition, the same rules apply for related words which are provided below.} \\
\\ 
\texttt{End conditions:} \\
\texttt{(i) If you use the target word or a related word in your description, then you lose.} \\
\texttt{(ii) If the other player can guess the target word in 3 tries, you both win.} \\
\\ 
\texttt{Let us start.} \\
\\ 
\texttt{This is the target word that you need to describe and that the other player needs to guess:} \\
\\ 
\texttt{guard} \\
\\ 
\texttt{Related words are:} \\
\\ 
\texttt{{-} security} \\
\texttt{{-} protect} \\
\texttt{{-} defend} \\
\\ 
\texttt{Important: You are under time pressure, give short descriptions that are to the point!} \\
            }
        }
    }
    & \\ \\

    \theutterance \stepcounter{utterance}  
    & \multicolumn{4}{p{0.6\linewidth}}{
        \cellcolor[rgb]{0.9,0.9,0.9}{
            \makecell[{{p{\linewidth}}}]{
                \texttt{\tiny{[P1$\rangle$GM]}}
                \texttt{<reasoning>} \\
\texttt{I need to come up with a clue for the word "guard" without using the word itself or any related words like security, protect, or defend. Let me think about what "guard" means. It's someone who protects a place or person, maybe like a guard dog or a security guard. But I can't use "security" or "protect" either.} \\
\\ 
\texttt{Hmm, maybe think about the action. A guard stands watch, keeps watch. So maybe something about watching or vigilance. Words like watch, vigil, lookout. But I need to make sure those aren't related. Wait, "watch" is a bit related to guard, but maybe it's allowed. Let me check the rules again. The related words are security, protect, defend. So "watch" might be okay. Alternatively, think of objects associated with guarding. Like a gate, a tower, a barrier. Or maybe something like a sentinel. But "sentinel" might be too close. } \\
\\ 
\texttt{Wait, the clue has to be a short description. Maybe use a metaphor. Like "keeps an eye on things" but "eye" is part of the body. Maybe "stands watch silently" but "watch" is okay? Or "prevents entry without permission" but that uses "prevent" which isn't a related word. Wait, the rules say I can't use the target word or parts/morphological variants, and also the related words provided. So "security" is a related word, so I can't use that. "Protect" and "defend" are also out. } \\
\\ 
\texttt{What about using a synonym for guard that's not in the related words? Like "warden" but that's a type of guard. Maybe "sentry"? But sentry is similar. Maybe "lookout"? That's someone who watches for danger. But "lookout" might be okay. Let me see. The clue needs to be a single phrase. Maybe "Watches for intruders silently" but "watches" is allowed? Or "Keeps an eye on a place" but "eye" is part of the body. Maybe "Acts as a barrier between safe and unsafe areas." Hmm, but "barrier" might be too vague. } \\
\\ 
\texttt{Alternatively, think of the tools a guard uses. A guard might have a flashlight, a baton, a badge. But those might not be specific enough. Or maybe the role} \\
            }
        }
    }
    & & \\ \\

    \theutterance \stepcounter{utterance}  
    & & & \multicolumn{2}{p{0.3\linewidth}}{
        \cellcolor[rgb]{0.9,0.9,0.9}{
            \makecell[{{p{\linewidth}}}]{
                \texttt{\tiny{[GM$|$GM]}}
                \texttt{abort game} \\
            }
        }
    }
    & & \\ \\

\end{supertabular}
}

\end{document}
